% ===========================================================================
\chapter{Camada de Serviços e Regras de Negócio (Service Layer)}
\label{cap_camada_servicos}
% ===========================================================================

A camada de serviços (\textit{Service Layer}) constitui a fronteira executiva da aplicação, na qual convergem a coordenação de operações de negócio, o controle de transações ACID e a imposição de regras invariantes institucionais e clínicas. Projetada sob o princípio da Inversão de Dependências (DIP) e gerenciada pelo container de Inversão de Controle (IoC) do Spring Framework, a camada opera de forma totalmente desacoplada de protocolos web ou mecanismos físicos de banco de dados.

% ===========================================================================
\section{Diagrama de Classes da Camada de Serviços}
\label{sec_diagrama_classes_servicos}
% ===========================================================================

A Figura \ref{fig_diagrama_classes_servicos} detalha as principais classes de serviço do SIGEA-GTT, suas dependências injetadas e as operações de negócio disponibilizadas aos clientes:

\begin{figure}[!ht]
\centering
\caption{Diagrama de Classes da Camada de Serviços (Spring Services)}
\label{fig_diagrama_classes_servicos}
\resizebox{0.96\textwidth}{!}{%
\begin{tikzpicture}[
    service_box/.style={
        rectangle split,
        rectangle split parts=3,
        rectangle split part fill={clinical-100!30, white, white},
        draw=clinical-700,
        thick,
        rounded corners=2pt,
        font=\tiny,
        minimum width=5.2cm,
        text width=5.0cm,
        align=left
    },
    dep/.style={
        draw=clinical-900,
        thick,
        dashed,
        -Latex
    }
]

    % ====================================================
    % AUTH SERVICE
    % ====================================================
    \node[service_box] (auth_serv) at (-4.0, 5.0) {
        \centering\textbf{\scriptsize <<Service>>\\AuthService}
        \nodepart{two}
        - userRepository : UserRepository\\
        - jwtTokenService : JwtTokenService\\
        - passwordEncoder : PasswordEncoder
        \nodepart{three}
        + authenticate(LoginRequestDTO) : LoginResponseDTO\\
        + firstLoginChangePassword(FirstLoginDTO) : void\\
        + changePassword(UUID, ChangePasswordDTO) : void
    };

    % ====================================================
    % JWT TOKEN SERVICE
    % ====================================================
    \node[service_box] (jwt_serv) at (4.0, 5.0) {
        \centering\textbf{\scriptsize <<Service>>\\JwtTokenService}
        \nodepart{two}
        - secretKey : String\\
        - expirationMs : long
        \nodepart{three}
        + generateToken(User) : String\\
        + extractEmail(String) : String\\
        + extractRole(String) : String\\
        + isTokenValid(String, User) : boolean
    };

    % ====================================================
    % USER SERVICE
    % ====================================================
    \node[service_box] (user_serv) at (-4.0, 0.0) {
        \centering\textbf{\scriptsize <<Service>>\\UserService}
        \nodepart{two}
        - userRepository : UserRepository\\
        - userMapper : UserMapper\\
        - passwordEncoder : PasswordEncoder
        \nodepart{three}
        + createUser(UserCreateDTO) : UserCreateResponseDTO\\
        + updateUser(UUID, UserUpdateDTO) : UserResponseDTO\\
        + updateStatus(UUID, UserStatusDTO) : void\\
        + getUsers(Pageable, UserRole) : Page<UserResponseDTO>\\
        + getUserById(UUID) : UserResponseDTO
    };

    % ====================================================
    % ACADEMIC CLASS SERVICE
    % ====================================================
    \node[service_box] (class_serv) at (4.0, 0.0) {
        \centering\textbf{\scriptsize <<Service>>\\AcademicClassService}
        \nodepart{two}
        - classRepository : AcademicClassRepository\\
        - userRepository : UserRepository\\
        - submissionRepository : ActivitySubmissionRepository
        \nodepart{three}
        + createClass(ClassCreateDTO, UUID) : ClassResponseDTO\\
        + enrollStudents(UUID, Set<UUID>) : void\\
        + removeStudent(UUID, UUID) : void\\
        + closeClass(UUID) : void\\
        + getClassesByProfessor(UUID, Pageable) : Page<ClassResponseDTO>
    };

    % ====================================================
    % ACTIVITY SERVICE
    % ====================================================
    \node[service_box] (act_serv) at (-4.0, -5.5) {
        \centering\textbf{\scriptsize <<Service>>\\ActivityService}
        \nodepart{two}
        - activityRepository : ActivityRepository\\
        - classRepository : AcademicClassRepository
        \nodepart{three}
        + createActivity(UUID, ActivityCreateDTO) : ActivityResponseDTO\\
        + getActivityDetail(UUID) : ActivityDetailDTO\\
        + getActivitiesByClass(UUID) : List<ActivityResponseDTO>
    };

    % ====================================================
    % ACTIVITY SUBMISSION SERVICE
    % ====================================================
    \node[service_box] (subm_serv) at (4.0, -5.5) {
        \centering\textbf{\scriptsize <<Service>>\\ActivitySubmissionService}
        \nodepart{two}
        - submissionRepository : ActivitySubmissionRepository\\
        - activityRepository : ActivityRepository\\
        - userRepository : UserRepository
        \nodepart{three}
        + submitActivity(UUID, UUID, SubmissionCreateDTO) : SubmissionResponseDTO\\
        + gradeSubmission(UUID, SubmissionGradeDTO) : SubmissionResponseDTO\\
        + getSubmissionsByActivity(UUID) : List<SubmissionResponseDTO>\\
        + getStudentSubmission(UUID, UUID) : Optional<SubmissionResponseDTO>
    };

    % Dependências
    \draw[dep] (auth_serv) -- node[above, font=\tiny] {usa} (jwt_serv);
    \draw[dep] (class_serv) -- (subm_serv);

\end{tikzpicture}
}
\fonte{Elaborado pelo autor (2026).}
\end{figure}

% ===========================================================================
\section{Detalhamento das Regras de Negócio Invariantes}
\label{sec_detalhamento_regras_servicos}
% ===========================================================================

Os serviços de negócio atuam como guardiões da consistência do sistema, validando fluxos antes de qualquer persistência:

\begin{enumerate}
    \item \textbf{Gestão de Identidade e Segurança (\texttt{AuthService} e \texttt{UserService})}:
    \begin{itemize}
        \item \textit{Validação de Domínio Institucional}: O método de criação de usuário intercepta o e-mail e aplica o validador \texttt{@UfsEmail}, rejeitando qualquer sufixo diferente de \texttt{@academico.ufs.br};
        \item \textit{Segregação de Matrícula Discente}: O serviço checa se o perfil do usuário é discente (\texttt{STUDENT}). Se for discente, a matrícula é obrigatória e deve possuir entre 6 e 30 caracteres alfanuméricos. Se for \texttt{ADMIN} ou \texttt{PROFESSOR}, qualquer tentativa de informar matrícula é forçadamente anulada;
        \item \textit{Fluxo de Senha Provisória}: No primeiro acesso, o método de login verifica se \texttt{mustChangePassword == true}. O usuário é impedido de acessar endpoints protegidos até que invoque com sucesso \texttt{firstLoginChangePassword}, que atualiza a credencial com BCrypt e define a flag provisória como falsa.
    \end{itemize}

    \item \textbf{Governança Acadêmica de Turmas (\texttt{AcademicClassService})}:
    \begin{itemize}
        \item \textit{Padrão Canônico de Nomenclatura}: O serviço concatena e valida a máscara textual universal \texttt{"Disciplina - Turma - Período"} (ex.: \textit{"Enfermagem na Saúde do Adulto I - T01 - 2025.2"});
        \item \textit{Docência Titular Obrigatória}: Toda turma é vinculada ao ID do professor autenticado, impedindo criação de turmas anônimas ou orfãs;
        \item \textit{Bloqueio de Encerramento com Pendências}: Antes de atualizar \texttt{isClosed = true}, o serviço invoca \texttt{countByActivityAcademicClassIdAndGradeIsNull(classId)}. Caso existam atividades com submissões não corrigidas pelo professor, lança-se uma \texttt{BusinessException} impedindo o encerramento.
    \end{itemize}

    \item \textbf{Auditoria Clínica e Avaliação Docente (\texttt{ActivitySubmissionService})}:
    \begin{itemize}
        \item \textit{Verificação de Janela de Entrega}: O serviço verifica se a data/hora atual é estritamente anterior a \texttt{activity.getDeadline()};
        \item \textit{Submissão Única}: A tentativa de reenvio por discente que já possua registro na atividade é bloqueada;
        \item \textit{Avaliação com Nota Numérica}: O método de correção exige que a nota esteja no intervalo $[0{,}00, 10{,}00]$ e que o parecer pedagógico (\texttt{professorFeedback}) contenha fundamentação formativa.
    \end{itemize}

    \item \textbf{Cálculo de Indicadores Clínicos IHI (\texttt{ClassDashboardService})}:
    \begin{itemize}
        \item Consolida as métricas recomendadas pelo IHI \cite{ihi2009}: Taxa de Eventos por 1.000 pacientes-dia, Taxa por 100 admissões e Percentual de internações com dano real (\% de categorias E a I).
    \end{itemize}
\end{enumerate}
